\documentclass[11pt]{article}
\usepackage[margin=1in]{geometry}
\usepackage{amsmath,amssymb,amsthm}

\newcommand{\Form}{\mathrm{Form}}
\newcommand{\T}{\mathrm{T}}
\newcommand{\F}{\mathrm{F}}
\newcommand{\TF}{\{\T,\F\}}
\newcommand{\ext}[1]{\overline{#1}}
\newcommand{\Sfive}{\{\lnot,\land,\lor,\to,\leftrightarrow\}}

\theoremstyle{definition}
\newtheorem{problem}{Problem}

\begin{document}

\begin{center}
{\LARGE Problem Set: The Interpretation of Propositional Formulas}\\[4pt]
{\large Goldrei \S 2.3, with emphasis on Theorem 2.2}
\end{center}

\bigskip

\noindent\emph{Throughline.} Theorem 2.2 has two halves. Uniqueness says a truth
assignment has no freedom beyond its values on the propositional variables;
existence says that on the variables the freedom is total. Much of this sheet is
about knowing which half you are using at any given moment --- and recognizing the
situations in which neither is available.

\subsection*{Ground rules}

\begin{itemize}
\item Notation and conventions follow Goldrei, Chapter 2: formulas are fully
bracketed; the \emph{length} of a formula is the number of occurrences of
connectives in it; for a function $v\colon P\to\TF$, $\ext{v}$ denotes the unique
truth assignment extending $v$ given by Theorem 2.2. As in the book, ``all truth
assignments'' abbreviates ``all truth assignments defined on a set of
propositional variables including those appearing in the formulas at hand.''
\item You may use without proof: Theorem 2.2, anything established in the book up
to the end of \S 2.3, and any earlier problem on this sheet.
\item \textbf{Refutation discipline.} To refute a claim of the form ``for all
formulas $\varphi,\psi,\dots$'' you must produce a fully concrete instance:
(i)~name a specific formula for each metavariable; (ii)~verify that your formulas
satisfy \emph{every} hypothesis of the claim; (iii)~exhibit a specific truth
assignment, by listing its values on the relevant propositional variables, and
compute the truth values needed to show that the conclusion fails. An explanation
of why a counterexample ought to exist is not a counterexample.
\item Starred problems are harder.
\end{itemize}

\section*{Part I --- Reading Theorem 2.2}

\begin{problem}[Reading the theorem precisely]
Let $P$ be a set of propositional variables and $S = \Sfive$. For each statement
below, decide whether it is true or false. Justify each verdict briefly; for false
statements, follow the refutation discipline.
\begin{enumerate}
\item[(a)] For every function $v\colon P\to\TF$ there is \emph{at least one} truth
assignment $w$ with $w(p)=v(p)$ for all $p\in P$.
\item[(b)] For every function $v\colon P\to\TF$ there is \emph{at most one} truth
assignment $w$ with $w(p)=v(p)$ for all $p\in P$.
\item[(c)] For every function $v\colon P\to\TF$ there is exactly one
\emph{function} $w\colon\Form(P,S)\to\TF$ with $w(p)=v(p)$ for all $p\in P$.
\item[(d)] For every truth assignment $w\colon\Form(P,S)\to\TF$ there is exactly
one function $v\colon P\to\TF$ such that $w=\ext{v}$.
\item[(e)] If $u,w$ are truth assignments with $u(p)=w(p)$ for all $p\in P$, then
$u(\varphi)=w(\varphi)$ for all $\varphi\in\Form(P,S)$.
\item[(f)] If $u$ is a truth assignment, $w\colon\Form(P,S)\to\TF$ is a function,
and $u(p)=w(p)$ for all $p\in P$, then $w$ is a truth assignment.
\item[(g)] If $u,w$ are truth assignments with $u(\varphi)=w(\varphi)$ for every
formula $\varphi$ of length at least $1$, then $u=w$.
\end{enumerate}
\end{problem}

\begin{problem}[Extensions that are not truth assignments]
Let $P=\{p\}$, $S=\Sfive$, and let $v\colon P\to\TF$ be given by $v(p)=\T$.
\begin{enumerate}
\item[(a)] Exhibit two distinct functions $f,g\colon\Form(P,S)\to\TF$ with
$f(p)=g(p)=\T$, neither of which is a truth assignment. For each, name a specific
formula at which it violates a specific truth-table clause, and verify the
violation.
\item[(b)] Show that there are infinitely many functions
$\Form(P,S)\to\TF$ extending $v$.
\item[(c)] In light of (b), exactly how many truth assignments extend $v$? Which
result are you citing?
\item[(d)] Prove that restriction to $P$ and the map $v\mapsto\ext{v}$ are
mutually inverse bijections between the set of truth assignments on $\Form(P,S)$
and the set of functions $P\to\TF$. Deduce that if $|P|=n$ then there are exactly
$2^n$ truth assignments on $\Form(P,S)$ (cf.\ Exercise~2.19).
\end{enumerate}
\end{problem}

\begin{problem}[What unique readability is for]
Work with strings over the alphabet $\{p,q,r,\land,\lor\}$ and \emph{no brackets}.
Define the \emph{pseudo-formulas} to be the smallest class of strings such that
each of $p,q,r$ is a pseudo-formula, and whenever $\alpha,\beta$ are
pseudo-formulas, so are the strings $\alpha\land\beta$ and $\alpha\lor\beta$.
Consider the string $\sigma$ given by $p\land q\lor r$.
\begin{enumerate}
\item[(a)] Exhibit two distinct construction trees for $\sigma$.
\item[(b)] Let $v(p)=\F$, $v(q)=v(r)=\T$. Compute a truth value for $\sigma$
bottom-up along each of your two trees, using the usual truth tables for $\land$
and $\lor$. Conclude that the recursive clauses do \emph{not} determine a function
from pseudo-formulas to $\TF$.
\item[(c)] Identify the sentence in the existence half of Goldrei's proof of
Theorem 2.2 that fails for pseudo-formulas, and explain exactly what goes wrong
with it here.
\item[(d)] Consider the attempted definition ``$\ext{v}(\sigma) :=$ the value
computed from the decomposition of $\sigma$.'' Say precisely which of
\emph{existence}, \emph{uniqueness}, and \emph{well-definedness} of $\ext{v}$ is
threatened by (b), and explain how these three notions are related in this
situation.
\end{enumerate}
\end{problem}

\begin{problem}[Adding a connective]
Let $\oplus$ (``exclusive or'') be a new binary connective whose truth table is
\[
\begin{array}{cc|c}
v(\theta) & v(\psi) & v((\theta\oplus\psi))\\ \hline
\T & \T & \F\\
\T & \F & \T\\
\F & \T & \T\\
\F & \F & \F
\end{array}
\]
and let $S_{\oplus}=\{\lnot,\land,\lor,\to,\leftrightarrow,\oplus\}$.
\begin{enumerate}
\item[(a)] Write out the clause ``$v$ respects the truth table of $\oplus$'' in
the style of \S 2.3.
\item[(b)] State the analogue of Theorem 2.2 for the language
$\Form(P,S_{\oplus})$.
\item[(c)] Describe exactly what must be added or changed in each half of the
proof --- the existence half and the uniqueness half --- to prove your statement
from (b).
\item[(d)] The proof relies on facts about $\Form(P,S_{\oplus})$ of the kind
established for $\Form(P,S)$ in \S 2.2. Identify these facts precisely, and
explain what would go wrong in your argument from (c) if the uniqueness of the
principal connective failed for the enlarged language.
\end{enumerate}
\end{problem}

\section*{Part II --- Computing with truth assignments}

\begin{problem}[Bottom-up evaluation]
Let $w$ be a truth assignment with $w(p)=\F$, $w(q)=\T$, $w(r)=\T$.
\begin{enumerate}
\item[(a)] Compute $w(((\lnot p\lor q)\to(r\land\lnot q)))$, listing the value of
every subformula.
\item[(b)] Compute $w(((p\to(q\to r))\leftrightarrow\lnot(p\lor\lnot r)))$,
listing the value of every subformula.
\end{enumerate}
For (c) and (d), say that the information above \emph{determines} $w(\varphi)$ if
every truth assignment $w$ satisfying $w(p)=\F$, $w(q)=\T$, $w(r)=\T$ gives
$\varphi$ the same value. (Note that nothing has been said about $w(s)$.)
\begin{enumerate}
\item[(c)] Is $w(((q\land s)\lor r))$ determined? If so, compute it and justify;
if not, exhibit two truth assignments consistent with the information above that
give the formula different values.
\item[(d)] Same question for $w(((q\land s)\land r))$.
\item[(e)] Now suppose all you know is $w(p)=\T$. For each of
$(p\lor q)$, $(p\land q)$, $(q\to p)$, $(q\lor\lnot q)$, $(q\land\lnot q)$,
decide whether $w(\varphi)$ is determined, and justify as in (c) and (d).
\end{enumerate}
\end{problem}

\begin{problem}[Truth tables, concrete and schematic]
\begin{enumerate}
\item[(a)] Give the full truth table of $((p\lor\lnot q)\to(q\land r))$.
\item[(b)] Give the truth table of $((\varphi\to\psi)\to\varphi)$ in terms of
the truth values of the formulas $\varphi$ and $\psi$, as in Exercise 2.23.
\item[(c)] Realize the row $(\varphi,\psi)=(\F,\T)$ of your table concretely:
produce specific formulas $\varphi_0,\psi_0$ built from $p,q$ and a truth
assignment $w$ (give its values on $p$ and $q$) with $w(\varphi_0)=\F$ and
$w(\psi_0)=\T$, then confirm by computation that
$w(((\varphi_0\to\psi_0)\to\varphi_0))$ agrees with your table.
\item[(d)] Now take $\psi$ to be the formula $\lnot\varphi$. Which rows of your
table in (b) remain realizable, and why do the others drop out? Show that
$((\varphi\to\lnot\varphi)\to\varphi)$ is a tautology if and only if $\varphi$
is a tautology, and a contradiction if and only if $\varphi$ is a contradiction.
\item[(e)] Using truth tables or direct argument, classify each of the following
as a tautology, a contradiction, or neither:
\[
((p\to q)\lor(q\to p)),\qquad
(p\leftrightarrow\lnot p),\qquad
((p\lor q)\to p),\qquad
(\lnot(p\land q)\leftrightarrow(\lnot p\lor\lnot q)).
\]
\end{enumerate}
\end{problem}

\begin{problem}[Counting assignments]
Let $n\ge 3$ and work with truth assignments on $\{p_1,\dots,p_n\}$. How many of
the $2^n$ truth assignments make each of the following formulas true?
\begin{enumerate}
\item[(a)] $(p_1\land\lnot p_2)$
\item[(b)] $\lnot((p_1\lor p_2)\lor p_3)$
\item[(c)] $\lnot(p_1\to p_2)$
\item[(d)] $((p_1\lor p_2)\land p_3)$, in the case $n=3$.
\end{enumerate}
Justify each count.
\end{problem}

\section*{Part III --- Inductions on the length of a formula}

\begin{problem}[The relevance lemma]
\begin{enumerate}
\item[(a)] Prove: if $u,w$ are truth assignments and $\varphi\in\Form(P,S)$,
$S=\Sfive$, is a formula such that $u(q)=w(q)$ for every propositional variable
$q$ occurring in $\varphi$, then $u(\varphi)=w(\varphi)$. Use induction on the
length of $\varphi$; write the base case and the cases $\lnot$, $\land$, $\to$ in
full, declaring the remaining cases similar. Make your induction hypothesis
explicit, and flag the exact steps at which you use that $u$ and $w$ respect the
truth tables.
\item[(b)] Deduce the statement of Exercise 2.18, together with its analogue for
the full connective set $S$.
\item[(c)] Deduce that the truth-table method is sound: if the propositional
variables occurring in $\varphi$ are among $p_1,\dots,p_n$, then $\varphi$ is a
tautology if and only if every one of the $2^n$ truth assignments on
$\Form(\{p_1,\dots,p_n\},S)$ makes $\varphi$ true. Be careful: a truth assignment
on $\{p_1,\dots,p_n\}$ and a truth assignment on a larger set $P$ of variables are
functions with different domains, and part of your task is to bridge the two.
\end{enumerate}
\end{problem}

\begin{problem}[Negating the variables: Exercise 2.28 extended]
For a truth assignment $v$, let $v^*$ be the truth assignment determined, via
Theorem 2.2, by
\[
v^*(p)=\begin{cases}\T,&\text{if }v(p)=\F,\\[2pt] \F,&\text{if }v(p)=\T,\end{cases}
\qquad\text{for all propositional variables }p.
\]
For $\varphi\in\Form(P,S)$, $S=\Sfive$, define $\varphi^*$ by recursion:
$p^*=\lnot p$ for each variable $p$; $(\lnot\theta)^*=\lnot\theta^*$; and
$(\theta\circ\psi)^*=(\theta^*\circ\psi^*)$ for each
$\circ\in\{\land,\lor,\to,\leftrightarrow\}$. (As always, this definition is
unambiguous thanks to unique readability.)
\begin{enumerate}
\item[(a)] Prove that $v(\varphi)=v^*(\varphi^*)$ for every truth assignment $v$
and every formula $\varphi$, by induction on the length of $\varphi$. Write the
base case and the cases $\lnot$, $\to$, $\leftrightarrow$ in full. (Exercise
2.28(a) is the restriction of this statement to $\{\lnot,\land,\lor\}$.)
\item[(b)] Show that $(v^*)^*=v$ \emph{as truth assignments}, being explicit about
where Theorem 2.2 is used. Deduce that $v\mapsto v^*$ is a bijection from the set
of truth assignments to itself.
\item[(c)] Deduce that $\varphi$ is a tautology if and only if $\varphi^*$ is a
tautology. Write the quantifier argument out in full: given an arbitrary truth
assignment $w$, say exactly which assignment you feed into (a), and why every
truth assignment arises in the required form.
\item[(d)] Deduce, without any new induction, that $\varphi$ is a contradiction if
and only if $\varphi^*$ is a contradiction.
\item[(e)] Prove or refute, without any new induction:
$v(\varphi^*)=v^*(\varphi)$ for every truth assignment $v$ and every formula
$\varphi$.
\end{enumerate}
\end{problem}

\begin{problem}[$\star$ The substitution lemma]
Fix a propositional variable $p$ and a formula $\chi$. For each formula
$\varphi\in\Form(P,S)$, $S=\Sfive$, define the formula $\varphi[\chi/p]$
(``$\varphi$ with $\chi$ substituted for $p$'') by recursion:
\[
p[\chi/p]=\chi;\qquad
q[\chi/p]=q\ \text{ for each variable }q\ne p;\qquad
(\lnot\theta)[\chi/p]=\lnot(\theta[\chi/p]);
\]
\[
(\theta\circ\psi)[\chi/p]=(\theta[\chi/p]\circ\psi[\chi/p])
\quad\text{for each }\circ\in\{\land,\lor,\to,\leftrightarrow\}.
\]
\begin{enumerate}
\item[(a)] Explain in one or two sentences why this recursion specifies
$\varphi[\chi/p]$ unambiguously, citing the relevant fact from \S 2.2.
\item[(b)] Prove: for every truth assignment $v$,
\[
v(\varphi[\chi/p])=u(\varphi),
\]
where $u$ is the truth assignment determined by $u(p)=v(\chi)$ and $u(q)=v(q)$ for
every variable $q\ne p$. (Which result guarantees that such a $u$ exists and is
unique?) Use induction on the length of $\varphi$; write the base case --- both
subcases --- and the cases $\lnot$ and $\to$ in full.
\item[(c)] Deduce: if $\varphi$ is a tautology, then $\varphi[\chi/p]$ is a
tautology, for every formula $\chi$.
\item[(d)] Show that the converse of (c) fails, following the refutation
discipline.
\item[(e)] In at most three sentences, use (b) to explain the following
phenomenon: in a schematic truth table such as the one in Problem 6(b), every row
is realizable when $\varphi$ and $\psi$ are distinct propositional variables, yet
rows can become unrealizable once $\varphi$ and $\psi$ are correlated, as in
Problem 6(d).
\end{enumerate}
\end{problem}

\section*{Part IV --- Tautologies, contradictions, counterexamples}

\begin{problem}[Calibration: true or false?]
For each claim, decide whether it is true or false. Give a short proof for each
true claim; follow the refutation discipline for each false one. Here $u,w,v$
range over truth assignments and $\varphi,\psi$ over formulas.
\begin{enumerate}
\item[(a)] If $u(\varphi)=w(\varphi)$ for some formula $\varphi$, then
$u(q)=w(q)$ for every propositional variable $q$ occurring in $\varphi$.
\item[(b)] There is a truth assignment $v$ with $v((p\land q))=\T$ and
$v(q)=\F$.
\item[(c)] There is a function $f\colon\Form(P,S)\to\TF$ with $f((p\land q))=\T$
and $f(q)=\F$.
\item[(d)] If $\varphi$ is not a tautology, then $\lnot\varphi$ is a tautology.
\item[(e)] If $\varphi$ is not a tautology, then $\lnot\varphi$ is not a
contradiction.
\item[(f)] If neither $\varphi$ nor $\psi$ is a tautology, then
$(\varphi\lor\psi)$ is not a tautology.
\item[(g)] If $(\varphi\land\psi)$ is a tautology, then $\varphi$ is a tautology.
\item[(h)] If $\varphi$ is not a contradiction, then some truth assignment makes
$\varphi$ true.
\end{enumerate}
\end{problem}

\begin{problem}[Around modus ponens: Exercise 2.27 revisited]
Let $\varphi,\psi$ be formulas. Prove or refute each of the following;
refutations must follow the refutation discipline.
\begin{enumerate}
\item[(a)] If $\varphi$ and $\psi$ are tautologies, then $(\varphi\land\psi)$ is a
tautology.
\item[(b)] If $(\varphi\land\psi)$ is a tautology, then $\varphi$ and $\psi$ are
both tautologies.
\item[(c)] If $(\varphi\lor\psi)$ is a tautology, then $\varphi$ is a tautology or
$\psi$ is a tautology.
\item[(d)] If $(\varphi\to\psi)$ is a tautology and $\psi$ is a contradiction,
then $\varphi$ is a contradiction.
\item[(e)] If $(\varphi\leftrightarrow\psi)$ is a tautology and $\varphi$ is a
tautology, then $\psi$ is a tautology.
\item[(f)] If $(\varphi\leftrightarrow\psi)$ is a tautology, then $\varphi$ and
$\psi$ are tautologies.
\end{enumerate}
\end{problem}

\begin{problem}[Forced, impossible, or independent]
An \emph{instance} of a hypothesis $H$ about metavariables $\varphi,\psi$ is a
choice of actual formulas for the metavariables under which $H$ holds. Given $H$
and a conclusion $C$, exactly one of the following obtains:
\begin{enumerate}
\item[(I)] every instance of $H$ satisfies $C$ \quad (``$H$ forces $C$'');
\item[(II)] no instance of $H$ satisfies $C$ \quad (``$H$ forces the failure of
$C$'');
\item[(III)] some instance of $H$ satisfies $C$ and some instance of $H$ does not
\quad (``$C$ is independent of $H$'').
\end{enumerate}
Note the difference between (II) and the mere failure of (I): in case (III) the
conclusion is not forced --- it can go either way --- whereas in case (II) it is
forced to fail.

For each pair below, decide which of (I), (II), (III) obtains. In cases (I) and
(II), prove the forcing. In case (III), give \emph{two} fully concrete instances
--- one where $C$ holds and one where it fails --- verifying in each instance
every part of $H$ and the status of $C$, per the refutation discipline.
\begin{enumerate}
\item[(a)] $H$: $\psi$ is a tautology. \quad $C$: $(\varphi\to\psi)$ is a
tautology.
\item[(b)] $H$: $(\varphi\to\psi)$ and $\psi$ are tautologies. \quad $C$:
$\varphi$ is a tautology.
\item[(c)] $H$: $\varphi$ is a contradiction. \quad $C$: $(\varphi\to\psi)$ is a
tautology.
\item[(d)] $H$: $\varphi$ is a contradiction. \quad $C$: $(\varphi\land\psi)$ is a
tautology.
\item[(e)] $H$: $(\varphi\to\psi)$ is a contradiction. \quad $C$: $\psi$ is a
tautology.
\item[(f)] $H$: $\varphi$ is not a tautology. \quad $C$: $\lnot\varphi$ is a
tautology.
\item[(g)] Parts (d) and (e) tacitly use the fact that no formula is both a
tautology and a contradiction. Prove this fact carefully. Which half of Theorem
2.2 --- existence or uniqueness --- does your proof use?
\end{enumerate}
\end{problem}

\begin{problem}[Characterizations: Exercise 2.31 revisited]
Let $\varphi,\psi$ be formulas.
\begin{enumerate}
\item[(a)] Suppose $\varphi$ is a tautology. Under what circumstances, if any, is
$(\varphi\to\psi)$ a tautology? State and prove a characterization of the form
``\dots\ if and only if \dots''.
\item[(b)] Suppose $\psi$ is a tautology. Under what circumstances, if any, is
$(\varphi\to\psi)$ a contradiction? Prove your answer.
\item[(c)] If $(\varphi\land\psi)$ is a contradiction, must $\varphi$ or $\psi$ be
a contradiction? Prove or refute.
\item[(d)] Suppose $(\varphi\lor\psi)$ is a contradiction. State and prove the
strongest conclusion you can about $\varphi$ and $\psi$, in the form of a
characterization of when $(\varphi\lor\psi)$ is a contradiction.
\item[(e)] If $(\varphi\leftrightarrow\psi)$ is a contradiction, must $\varphi$ or
$\psi$ be a contradiction? Prove or refute. Then state and prove a
characterization: $(\varphi\leftrightarrow\psi)$ is a contradiction if and only if
\dots\ (a condition relating the values of $\varphi$ and $\psi$ under each truth
assignment).
\end{enumerate}
\end{problem}

\begin{problem}[Majority rule: Exercise 2.29 revisited]
Let $S$ be a finite set of truth assignments with $|S|$ odd, and let $D$ be the
set of formulas made true by more than half of the members of $S$.
\begin{enumerate}
\item[(a)] Show that for every formula $\varphi$, exactly one of $\varphi$,
$\lnot\varphi$ belongs to $D$.
\item[(b)] Show that if $\varphi\in D$ and $(\varphi\to\theta)$ is a tautology,
then $\theta\in D$.
\item[(c)] Construct a set $S$ of three truth assignments on $\{p,q\}$ ---
specify each assignment by its values on $p$ and $q$ --- such that $p\in D$ and
$(p\to q)\in D$ but $q\notin D$. Verify all three memberships.
\item[(d)] If $\varphi\in D$ and $\psi\in D$, must $(\varphi\land\psi)\in D$?
Prove, or refute with a concrete $S$ and concrete formulas.
\item[(e)] Now drop the requirement that $|S|$ is odd: exhibit a set $S$ with
$|S|=2$ for which the conclusion of (a) fails.
\item[(f)] Show that if $|S|=1$ then $D$ is closed under modus ponens: if
$\varphi\in D$ and $(\varphi\to\theta)\in D$, then $\theta\in D$. Identify the
precise step of your argument that breaks down when $|S|=3$, and explain how your
example in (c) exploits exactly that gap.
\end{enumerate}
\end{problem}

\section*{Part V --- Capstone}

\begin{problem}[$\star$ Repetition-free formulas: Exercise 2.30 generalized]
Call a formula \emph{repetition-free} if no propositional variable occurs more
than once in it.
\begin{enumerate}
\item[(a)] \emph{(Gluing Lemma.)} Let $\theta$ and $\psi$ be formulas with no
propositional variable in common, and let $u,w$ be truth assignments. Show that
there is a truth assignment $x$ with $x(\theta)=u(\theta)$ and
$x(\psi)=w(\psi)$. Define $x$ explicitly on the propositional variables, and cite
Theorem 2.2 and Problem 8 at the precise points where each is used.
\item[(b)] Prove by induction on length: every repetition-free formula of
$\Form(P,S)$, $S=\Sfive$, is made true by some truth assignment and made false by
some truth assignment. A complete proof should record why the immediate
subformulas of a repetition-free formula are themselves repetition-free and, in
the binary cases, have no variable in common; it should also identify the feature
shared by the truth tables of all five connectives that makes the inductive step
go through.
\item[(c)] Deduce Exercise 2.30, and state the stronger conclusion that your proof
of (b) actually delivers.
\item[(d)] Exhibit a specific formula containing a repeated variable, and point to
the specific step of your inductive argument in (b) that fails for it.
\item[(e)] In one short paragraph: the freedom exploited in (a) is exactly what is
\emph{not} available in Problem 13(b) (equivalently, Exercise 2.27(c)). Explain
why distinct metavariables $\varphi,\psi$ need not behave like distinct
propositional variables, and which conditions on $\varphi$ and $\psi$ would
restore the analogy.
\end{enumerate}
\end{problem}

\end{document}